% page 66 du fac-simile (image p-065 du scan)
% bloc 160.0 x 224.7 mm ; marge gauche 8.0 mm
\pgc{-12.700mm}{0.000mm}{
\l{\cel{6}58}
\l{}
\l{\sou{baritono}. I. Muzik-instrumento antiqua, quaze viulo, qua}
\l{\cel{3}havas, sur sua mancho, dek-e-sis kordi de latuno quin}
\l{\cel{3}onu pinchis per la fingri, ultre sep tripo-kordi a qui}
\l{\cel{3}onu efikis per arketo. - II. Homo-voco, inter tenoro e}
\l{\cel{3}basvoco. - DEFIRS.}
\l{}
\l{\sou{barko}. Mikra kanoto, sen-masta, movebla per remili, e gene-}
\l{\cel{3}rale ne-uzata en aquo marala. - DEFIRS.}
\l{}
\l{\sou{barkarolo}. Kansono ritmoza di la gondolisti di la urbo Vene-}
\l{\cel{3}zia (Italia). - DEFIRS.}
\l{}
\l{\sou{barkaso}. Granda barko uzata por la desembarko di la pasajeri}
\l{\cel{3}e di la vari kande ne existas portui ube la navi povas}
\l{\cel{3}an-kayeskar. - DFIRS.}
\l{}
\l{\sou{barografo}.\cel{2}Barometro enrejistrera, konstruktita specale por}
\l{\cel{3}donar, sur grafiko kontinua, la kurvo di la altitudi}
\l{\cel{3}atingita da aernavigero, aviacisto. - DEFIRS.}
\l{}
\l{\sou{baroka}.\cel{2}Qua aspektas nereguloza o stranja (ex.: perlo, stilo).}
\l{\cel{3}- DEFIS.}
\l{}
\l{\sou{barometro}.\cel{2}Instrumento fizikala qua uzesas : I. por mezurar}
\l{\cel{3}la pezo di aero, per recipiento apertita, plena ye merkurio}
\l{\cel{3}qua, presata da atmosfero, su elevas plu o min alte en}
\l{\cel{3}tubo gradizita, aero-vakua ; II. por mezurar la altitudi,}
\l{\cel{3}pro ke la aero-kolono deskreskas altesale, e ke desdens-}
\l{\cel{3}eskas atmosfero segun ke onu acensas ; III. por prognozar}
\l{\cel{3}vetero segun la denseso-grado di atmosfero. - DEFIRS.}
\l{}
\l{\sou{barono}. En la hierarkio di la tituli di nobeleso (grantita}
\l{\cel{3}da la suvereno, od heredala), olta qua venas pos la vis-}
\l{\cel{3}komto ed ante la+chevaliere.- DEFIRS.}
\l{}
\l{\sou{baroskopo}.\cel{2}Aparato por konstatar ke korpo pozita en aero}
\l{\cel{3}perdas parto de sua pezo, egala a la pezo di la aero-volu-}
\l{\cel{3}mino quan olu diplasas. - DEFIRS.}
\l{}
\l{\sou{barsho}. (\sou{zool.}) Fisho qua tre similesas perko. - L.\cel{1}\sou{labrax}}
\l{\cel{3}\sou{lupus}. - DFI.}
\l{}
\l{\sou{barto}. Omna ek la lami korno-harda qui garnisas transverse}
\l{\cel{3}la maxilo di la baleno e di ula cetacei. - D.}
\l{}
\l{\sou{basa}. I. Qua ne atingas la nivelo ordinara. - II. Qua ne}
\l{\cel{3}atingas la nivelo di altra objekto. - F.}
\l{}
\l{\sou{basamento}. Infra parto (bazo), videbla, di edifico, qua jacas}
\l{\cel{3}sur la fundamento. - DEFIS.}
}
